\section{Environmental selection and mixed source weights}\label{fp:sec:positive}

\subsection{Positive moment orders for the same scalar model}
Let $E,W$ be independent stationary standard Ornstein--Uhlenbeck
processes with covariance $e^{-|u-v|/2}$. We use the same conditional survival as in \eqref{pr:defF}, now allowing all positive moment orders. For $0<\rho<1$, set
\begin{equation}\label{fp:eq:basic}
 \beta=\sqrt{\rho/(1-\rho)},\qquad
 q_t(E)=\PP_W\{W_s\ge-\beta E_s\ (0\le s\le t)\mid E\},
 \qquad S_t=-\log q_t(E).
\end{equation}
For every $p>0$ define
\begin{equation}\label{fp:eq:positive-spectrum}
 F_{p,\rho}(t)=\EE q_t^p,\qquad b_t(p,\rho)=-\log F_{p,\rho}(t),
 \qquad \chi(p,\rho)=\lim_{t\to\infty}b_t(p,\rho)/t.
\end{equation}
Theorem~\ref{fp:thm:joint-constant} proves existence and, locally
uniformly for positive $p$ and $0<\rho<1$,
\begin{equation}\label{fp:eq:constant}
 b_t(p,\rho)=t\chi(p,\rho)+O(1),\qquad t\ge1.
\end{equation}
For $0<p<1$ this is the scalar in \eqref{profile:chi}.
The ordinary probability from \eqref{pr:ordinaryOU} is
\begin{equation}\label{fp:eq:ordinary}
 a(t)=\EE q_t=\pi^{-1}\arcsin(e^{-t/2}),\qquad
 -\log a(t)=t/2+O(1).
\end{equation}
Every continuous environmental path is bounded on a compact interval;
a sufficiently high private OU tube gives $q_t(E)>0$. The initial
constraint gives $q_t(E)<1$. Thus $S_t$ is finite almost surely.
These moments describe the selection of the common environment.
Mallein and Mi\l o\'s~\cite[Theorem~1.6]{MalleinMilos2018} established
quenched and annealed OU persistence exponents for admissible initial
data, with a strict gap between them for a nonzero wall amplitude.
Here the stationary moment spectrum, its constant-factor estimates
and its derivatives will locate the clocks supporting the arithmetic
lower bound.


\subsection{A dimension-free curvature estimate}
\begin{lemma}[Gaussian convex-probability curvature]\label{fp:lem:curvature}
Let $C\subset\R^d$ be convex with positive Gaussian measure and let
\[
 A(x)=-\log\PP\{Z+\beta x\in C\},\qquad Z\sim N(0,I_d),\quad\beta>0.
\]
Then
\begin{equation}\label{fp:eq:curvature}
 0\preceq D^2A(x)\preceq\beta^2I_d,\qquad
 |\nabla A(x)|^2\le2\beta^2A(x).
\end{equation}
Under $d\nu_p\propto e^{-pA}d\gamma_d$, $p>0$,
\begin{equation}\label{fp:eq:variance}
 \Var_{\nu_p}(A)\le2\beta^2\EE_{\nu_p}A.
\end{equation}
\end{lemma}
\begin{proof}
Let $U_x$ have law $N(\beta x,I_d)$ conditioned to lie in $C$. Differentiating the Gaussian convolution gives
\[
 D^2A(x)=\beta^2\{I_d-\Cov(U_x)\}.
\]
The variance inequality of Brascamp and Lieb~\cite{BrascampLieb1976}
for a strongly log-concave measure gives $\Cov(U_x)\preceq I_d$.
For the hard-boundary approximation, a convex set of positive Gaussian
measure has nonempty interior, and replacing it by its closure changes
no Gaussian integral. Use the potentials
\[
 V_{N,x}(u)=\tfrac12|u-\beta x|^2+N\operatorname{dist}(u,C)^2.
\]
Smooth convex approximations to the squared distance allow the classical inequality with $D^2V_{N,x}\succeq I_d$. First remove the smoothing and then let $N\to\infty$; Gaussian domination passes the means and second moments to the conditioned law. This proves the covariance bound for the possibly unbounded hard convex set. Covariance is positive semidefinite, which proves the other bound. We use the standard form
\[
 \Var_\nu f\le\EE_\nu\langle\nabla f,(D^2V)^{-1}\nabla f\rangle,
 \qquad d\nu\propto e^{-V}dx;
\]
see~\cite[equation~(1.3)]{CarlenCorderoLieb2013}. Since $A\ge0$ and its Hessian is bounded above by $\beta^2I_d$,
\[
 0\le A(x-\beta^{-2}\nabla A(x))
 \le A(x)-\frac{|\nabla A(x)|^2}{2\beta^2}.
\]
Finally, $|x|^2/2+pA(x)$ has Hessian at least $I_d$. Apply Brascamp--Lieb once more and use the gradient estimate.
\end{proof}

\begin{theorem}[Regularity of the positive spectrum]\label{fp:thm:regularity}
For $p>0$ and $0<\rho<1$, $\chi(p,\rho)$ is nondecreasing and concave in $p$, and belongs to $C^{1,1}_{\mathrm{loc}}$ in that variable. It is jointly locally Lipschitz in $(p,\rho)$, and $\partial_p\chi$ is jointly continuous. Locally,
\begin{equation}\label{fp:eq:rho-holder}
 |\partial_p\chi(p,\rho)-\partial_p\chi(p,\rho')|
 \le C|\rho-\rho'|^{1/2}.
\end{equation}
Moreover, uniformly on compact positive parameter sets,
\begin{equation}\label{fp:eq:derivative-rate}
 \partial_pb_t(p,\rho)=t\partial_p\chi(p,\rho)+O(\sqrt t),
 \qquad t\ge1.
\end{equation}
\end{theorem}
\begin{proof}
On a finite grid, simultaneously whitening the two equal covariance matrices gives the cone $\Sigma^{-1/2}\R_+^d$ in Lemma~\ref{fp:lem:curvature}. With $b=-\log\EE e^{-pA}$,
\[
 b'=\EE_{\nu_p}A,\qquad
 -b''=\Var_{\nu_p}(A)\le2\beta^2b'.
\]
Concavity and $b(0)=0$ imply $0\le b'(p)\le b(p)/p$. For $p\le1$, $q^p\ge q$; for $p\ge1$, Jensen gives $\EE q^p\ge(\EE q)^p$. Hence
\begin{equation}\label{fp:eq:linear-b}
 0\le b(p)\le\max(1,p)[-\log a(t)]
 \le\max(1,p)(t/2+\log\pi).
\end{equation}
Thus $b/t$ has uniformly locally bounded first derivatives and uniformly Lipschitz first derivatives for $t\ge1$. For nested grids whose union is dense and includes both endpoints, the conditional grid probabilities decrease to $q_t(E)$ by path continuity. Their negative logarithms therefore increase to $S_t$. On each compact positive order interval, the functions $x^j e^{-px}$, $j=0,1,2$, are bounded uniformly in $x\ge0$. Dominated convergence consequently passes the unnormalized moments and their first two order derivatives to the path limit. The normalizers are positive, so the same is true of the derivatives of $b_t$. Equicontinuity, pointwise convergence and integration of the first derivatives prove the $C^{1,1}$ assertion. The constant-factor estimate \eqref{fp:eq:constant}, proved in Appendix~\ref{fp:app:joint}, and difference quotients of length $t^{-1/2}$ give \eqref{fp:eq:derivative-rate}.

For the dependence on $\rho$, the finite-grid semigroup calculation
in the proof of Proposition~\ref{pr:lipschitz} gives
\[
 \left|\partial_u\log\EE q_\rho^p\right|
 \le C\frac{\EE[q_\rho^p(1-\log q_\rho)]}{\EE q_\rho^p},
 \qquad u=-\tfrac12\log\rho.
\]
The calculation uses $|p(1-p)|$, so it applies to all compact positive
order intervals. By concavity, the ratio is at most $1+b_t(p,\rho)/p$;
\eqref{fp:eq:linear-b} bounds it by $C(1+t)$. A slightly smaller
positive exponent dominates the logarithmic integrand when removing
smooth approximation and passing to dense grids. Integrating the
parameter bounds and taking the limiting slope proves joint local
Lipschitz continuity. Finally compare difference quotients at $\rho,\rho'$ with step $|\rho-\rho'|^{1/2}$, using the uniform Lipschitz bound for $\partial_p\chi$ in $p$. This proves \eqref{fp:eq:rho-holder} and joint continuity.
\end{proof}

\subsection{A neighboring-moment localization principle}
\begin{lemma}\label{fp:lem:localization}
Let $J\Subset(0,\infty)$ be a compact interval, let $\eta\in\operatorname{int}J$, and let $R\in[0,1]$, $\tau\ge0$. Suppose, for the same random variable $R$ at every $p\in J$,
\begin{equation}\label{fp:eq:moment-window}
 C^{-1}e^{-\tau\chi(p,\rho)}\le\EE R^p\le Ce^{-\tau\chi(p,\rho)}.
\end{equation}
Put $d\nu_R=R^\eta d\PP/\EE R^\eta$. Then
\begin{equation}\label{fp:eq:bernstein}
 \nu_R\{|{-\log R}-\tau\partial_\eta\chi(\eta,\rho)|>u\}
 \le C_1\exp\left[-c_1\min\left\{\frac{u^2}{1+\tau},u\right\}\right],
 \qquad u\ge0.
\end{equation}
The constants are locally uniform in the indicated parameters. Zeros of $R$ have zero $\nu_R$-mass.
\end{lemma}
\begin{proof}
For $|h|$ smaller than a fixed distance to $\partial J$,
\[
 \EE_{\nu_R}e^{h(-\log R-\tau\partial_\eta\chi)}
 =e^{-h\tau\partial_\eta\chi}\frac{\EE R^{\eta-h}}{\EE R^\eta}
 \le C^2e^{C_2\tau h^2},
\]
by the $C^{1,1}$ Taylor bound of Theorem~\ref{fp:thm:regularity}. Apply Chernoff in both directions with $|h|$ comparable to $\min\{u/(1+\tau),1\}$. 
\end{proof}

\subsection{A strictly positive cost of annealing first}
\begin{theorem}[Strict extensive annealing loss]\label{fp:thm:annealing}
For every $0<\eta,\rho<1$,
\begin{equation}\label{fp:eq:annealing-loss}
 \chi(\eta,\rho)-\frac\eta2
 \ge\frac{\rho}{16\pi}\min\{\eta,1-\eta\}>0.
\end{equation}
\end{theorem}
\begin{proof}
Let $\mathsf P_t$ be the environmental OU law and $\mathsf Q_t$ its law conditional on joint survival. Then
\[
 f_t:=\frac{d\mathsf Q_t}{d\mathsf P_t}=\frac{q_t}{a(t)}.
\]
For $H_t=\int_0^tE_s\,ds$, the unconditioned variance is
\[
 v_t=4t-8(1-e^{-t/2})\le4t.
\]
Complete $Z=\sqrt\rho E+\sqrt{1-\rho}W$ to an independent orthogonal Gaussian rotation. Only $Z$ is conditioned to stay positive. At time $s$ its density on $(0,\infty)$ is proportional to
\[
 \varphi(z)h_s(z)h_{t-s}(z),\qquad h_u(z)=\PP_z(\tau_0>u).
\]
Both factors are increasing, so the conditional law dominates a positive half-normal. Hence
\[
 \EE_{\mathsf Q_t}H_t\ge\sqrt{2\rho/\pi}\,t=:m_0t.
\]
Under $\mathsf Q_t$, $H_t$ is sub-Gaussian about its mean with variance proxy $v_t$. To verify this, append $H_t$ to a finite Gaussian survival grid. After any linear tilt, the survival constraint is convex, and Brascamp--Lieb bounds the variance of $H_t$ by its original Gaussian variance. Integrating the second derivative of the log moment generating function proves the claim. The finite-grid normalizers decrease to the positive number $a(t)$; Gaussian integrability passes the means and exponential moments to the limit.

For $B_t=\{H_t\ge m_0t/2\}$, the two sub-Gaussian estimates give
\[
 \mathsf P_t(B_t)\le e^{-\rho t/(16\pi)},\qquad
 \mathsf Q_t(B_t^c)\le e^{-\rho t/(16\pi)}.
\]
H\"older on $B_t$ and $B_t^c$ yields
\begin{align*}
 \int f_t^\eta\,d\mathsf P_t
 &\le\mathsf Q_t(B_t)^\eta\mathsf P_t(B_t)^{1-\eta}
   +\mathsf Q_t(B_t^c)^\eta\mathsf P_t(B_t^c)^{1-\eta}\\
 &\le2e^{-\rho t\min(\eta,1-\eta)/(16\pi)}.
\end{align*}
Use $F_{\eta,\rho}(t)=a(t)^\eta\int f_t^\eta d\mathsf P_t$, take negative logarithms, divide by $t$, and let $t\to\infty$.
\end{proof}

\begin{remark}[Exact information identity]\label{fp:rem:renyi}
With the standard R\'enyi divergence~\cite{VanErvenHarremoes2014},
\begin{equation}\label{fp:eq:renyi}
 -\log\EE q_t^\eta
 =-\eta\log a(t)+(1-\eta)D_\eta(\mathsf Q_t\Vert\mathsf P_t).
\end{equation}
Thus $\chi-\eta/2$ is a nonzero information rate for the change in environmental law caused by survival. Theorem~\ref{fp:thm:annealing} gives a positive lower bound for this rate.
\end{remark}


\subsection{Mixed weights and Gibbs variational error}\label{fp:sec:mixed}
We now consider the weight that arises when an accepted private path is divided by an environmental envelope. Let $J=[p_-,p_+]\Subset(0,1)$, $\eta\in\operatorname{int}J$, and
\[
 0\le q\le1,\qquad 0<D\le1,\qquad q\le C_0D,\quad C_0\ge1.
\]
For fixed $\tau\ge0$ and $0<\rho<1$, the same $q,D$ satisfy,
for every $p\in J$,
\begin{equation}\label{fp:eq:two-windows}
 \EE q^p\asymp\EE D^p\asymp e^{-\tau\chi(p,\rho)}.
\end{equation}
Set
\begin{equation}\label{fp:eq:mixed-measure}
 R=\frac qD,\quad w=RD^\eta=qD^{\eta-1},\quad
 Z_\eta=\EE w,\quad d\widehat\nu_\eta=\frac w{Z_\eta}\,d\PP.
\end{equation}
On the support of this measure write $S=-\log q$, $T=-\log D$. On $\{q=0\}$ all integrands with a strictly positive power of $q$ are defined as zero.

\begin{lemma}[Mixed-power stability]\label{fp:lem:mixed-power}
Uniformly for $p\in J$ and $a$ in a compact subset of $(0,\infty)$,
\begin{equation}\label{fp:eq:mixed-power}
 \EE[R^aD^p]\asymp e^{-\tau\chi(p,\rho)}.
\end{equation}
\end{lemma}
\begin{proof}
Put $A_p=\EE D^p$, $d\nu_p=D^p d\PP/A_p$ and $r=R/C_0\in[0,1]$. Then
\[
 \EE_{\nu_p}r^p
 =\frac{\EE q^p}{C_0^p\EE D^p}\ge\kappa>0.
\]
For $a\le p$, $r^a\ge r^p$; for $a\ge p$, Jensen gives $\EE r^a\ge(\EE r^p)^{a/p}$. Hence
\[
 \kappa^{\max(1,a/p)}\le\EE_{\nu_p}r^a\le1.
\]
Multiplication by $C_0^aA_p$ proves the result. In fact this lemma alone is valid on any compact positive $p$-interval.
\end{proof}

\begin{theorem}[Mixed spectrum and localization]\label{fp:thm:mixed}
The normalization satisfies $Z_\eta\asymp e^{-\tau\chi(\eta,\rho)}$. If
\begin{equation}\label{fp:eq:signed-domain}
 s<1,\qquad p:=\eta-s-t\in\operatorname{int}J,
\end{equation}
then, uniformly on compact subsets of this region,
\begin{equation}\label{fp:eq:mixed-mgf}
 \log\EE_{\widehat\nu_\eta}e^{sS+tT}
 =\tau[\chi(\eta,\rho)-\chi(\eta-s-t,\rho)]+O(1).
\end{equation}
Both $S$ and $T$ satisfy \eqref{fp:eq:bernstein} under $\widehat\nu_\eta$. In addition,
\begin{equation}\label{fp:eq:ratio-tail}
 \widehat\nu_\eta\{\log(D/q)>v\}\le Ce^{-v}\quad(v\ge0),
 \qquad \log(D/q)\ge-\log C_0,
\end{equation}
and
\begin{equation}\label{fp:eq:mixed-entropy}
 \KL(\widehat\nu_\eta\Vert\PP)
 =\tau[\chi(\eta,\rho)-\eta\chieta(\eta,\rho)]
 +O(\sqrt{1+\tau}).
\end{equation}
\end{theorem}
\begin{proof}
The identity
\[
 RD^\eta e^{sS+tT}=R^{1-s}D^{\eta-s-t}
\]
and Lemma~\ref{fp:lem:mixed-power} prove \eqref{fp:eq:mixed-mgf}. Neighboring values of $s$ or $t$ and the same Chernoff argument as in Lemma~\ref{fp:lem:localization} give concentration about $\tau\chieta$.

On $\{D/q>e^v\}$ one has $w\le e^{-v}D^\eta$. Divide expectations by $Z_\eta$ and use \eqref{fp:eq:two-windows}; this proves the tail bound. In particular $\EE_{\widehat\nu}|\log R|=O(1)$. Integrating the concentration bound gives $\EE_{\widehat\nu}T=\tau\chieta+O(\sqrt{1+\tau})$. Finally,
\[
 \log\frac{d\widehat\nu}{d\PP}=-\eta T+\log R-\log Z_\eta,
\]
which proves \eqref{fp:eq:mixed-entropy}.
\end{proof}

\begin{proposition}[Two Gibbs comparisons and near-minimality]\label{fp:prop:gibbs}
Let
\[
 d\nu_\eta^D=\frac{D^\eta}{\EE D^\eta}\,d\PP,
 \qquad
 d\nu_\eta^q=\frac{q^\eta}{\EE q^\eta}\,d\PP.
\]
Then
\begin{equation}\label{fp:eq:gibbs-density}
 \frac{d\widehat\nu_\eta}{d\nu_\eta^D}\le C,
 \qquad \frac{d\widehat\nu_\eta}{d\nu_\eta^q}\le C,
 \qquad
 \KL(\widehat\nu_\eta\Vert\nu_\eta^D),\,
 \KL(\widehat\nu_\eta\Vert\nu_\eta^q)\le\log C.
\end{equation}
Consequently the actual mixed measure has $O(1)$ excess in each of the two finite-time Gibbs variational problems:
\begin{align}
 0&\le\KL(\widehat\nu_\eta\Vert\PP)+\eta\EE_{\widehat\nu_\eta}T+\log\EE D^\eta\le\log C,\label{fp:eq:gibbs-D}\\
 0&\le\KL(\widehat\nu_\eta\Vert\PP)+\eta\EE_{\widehat\nu_\eta}S+\log\EE q^\eta\le\log C.\label{fp:eq:gibbs-q}
\end{align}
\end{proposition}
\begin{proof}
The two density ratios are respectively
\[
 R\frac{\EE D^\eta}{Z_\eta},\qquad
 R^{1-\eta}\frac{\EE q^\eta}{Z_\eta}.
\]
They are bounded because $R\le C_0$ and $0<\eta<1$. The normalization estimates give the claimed uniform constants. Relative entropy is nonnegative and is at most the logarithm of an upper density bound. Expand the two relative entropies; integrability was proved in Theorem~\ref{fp:thm:mixed}.
\end{proof}

\begin{remark}\label{fp:rem:forward-kl}
The bound relative to $\nu_\eta^q$ uses $\eta\le1$; the power-stability and mixed-moment formulas themselves have a larger positive-parameter range. Appendix~\ref{fp:sec:histories} treats finitely many conditional stages and groups of comparisons on a shared clock.
\end{remark}

The conditioned bridges and accepting measure constructed next provide
the fixed neighboring-moment windows needed to apply these estimates
to original prime tuples in Section~\ref{fp:sec:arithmetic}.
